% OUTDATED
\chapter{Formy normalne macierzy rzeczywistych}

Głównym celem pracy jest omówienie form normalnych macierzy zespolonych symetrycznych, antysymetrycznych oraz ortogonalnych.
Dlatego najpierw musimy omówić odpowiednie przypadki rzeczywiste macierzy symetrycznych, antysymetrycznych oraz ortogonalnych.
Formą normalną będziemy nazywać postać kanoniczną albo inaczej postać diagonalną macierzy lub najbardziej podobną postać do postaci diagonalnej.  
W tym rozdziale będą przedstawione podstawowe pojęcia oraz formy normalne dla macierzy rzeczywistych symetrycznych, antysymetrycznych oraz ortogonalnych.

\section{Podstawowe pojęcia}

Niech $V$ będzie przestrzenią liniową nad ciałem $\K$ liczb zespolonych lub rzeczywistych ($\K$ = $\C$ lub $\K$ = $\R$). Przekształcenie $\inner{\cdot}{\cdot} \colon V \times V \to \K$ przyporządkowujące parze wektorów ($x$, $y$) $\in$ $V \times V$ liczbę z $\K$ nazywamy \textbf{iloczynem skalarnym}, jeżeli dla dowolnych wektorów $x$, $y$, $z$ $\in$ $V$ spełnione są następujące warunki:

\begin{enumerate}[leftmargin=1cm, itemsep=2pt, label=(\arabic*)]
    \item $\inner{x}{y}$ = $\overline{\inner{y}{x}}$, gdzie $\overline{\inner{y}{x}}$ jest sprzężeniem zespolonym liczby $\inner{y}{x}$
    \item $\inner{ax}{y}$ = $a\inner{x}{y}$, gdzie $a \in \K$
    \item $\inner{x + y}{z}$ = $\inner{x}{z} + \inner{y}{z}$
    \item $\inner{x}{x} \geq 0$ oraz $\inner{x}{x} = 0 \iff x = 0$ (warunek o dodatniej określoności)
\end{enumerate}
Z (1), (2), (3) wynika, że:
\begin{enumerate}[leftmargin=1cm, itemsep=2pt, label=(\arabic*), start=5]
    \item $\inner{x}{ay}$ = $\overline{\inner{ay}{x}}$ = $\overline{a}\inner{x}{y}$, gdzie $a \in \K$
    \item $\inner{x}{y + z}$ = $\inner{x}{z} + \inner{x}{y}$
\end{enumerate}

\begin{definition}
Przestrzeń liniową $V$ nad ciałem liczb $\K$ wraz z wyróżnionym iloczynem skalarnym $\inner{\cdot}{\cdot}$, tj. parę ($V$, $\inner{\cdot}{\cdot}$) nazywamy $\textbf{przestrzenią unitarną}$.
W przypadku gdy $\K = \R$, przestrzeń unitarną nazywamy \textbf{przestrzenią euklidesową}.
\end{definition}

Dalej będziemy rozpatrywać tylko przestrzenie o skończonym wymiarze.

\begin{definition}
    \textbf{Normą wektora} $x \in V$ nazywamy liczbę $\norm{x} = \sqrt{\inner{x}{x}}$ 
\end{definition}

Niech $V$ będzie $n$-wymiarową przestrzenią unitarną z dowolną bazą $\base{e}$. Niech \textbf{x}, \textbf{y} będą dowolnymi wektorami z $V$, gdzie $x_i$, $y_i$ są współrzędnymi tych wektorów:

\begin{center}
    $$\textbf{x} = \sum_{i=1}^{n}{x_{i}e_{i}} \quad \textbf{y} = \sum_{i=1}^{n}{y_{i}e_{i}}$$
\end{center}
Wtedy z (2), (3), (5), (6):
\begin{center}
    $$\inner{\textbf{x}}{\textbf{y}} = \inner{
        \sum_{i=1}^{n}{x_{i} \textbf{e}_i}
    }{
        \sum_{k=1}^{n}{y_{k} \textbf{e}_k}
    } = \sum_{i=1}^{n}{\sum_{k=1}^{n}{x_{i} \bar{y}_{k} \inner{\textbf{e}_{i}}{\textbf{e}_{k}}}}$$
\end{center}
Niech
\begin{equation}\label{eq:def_h}
    h_{ik} = \inner{\textbf{e}_i}{\textbf{e}_k}, \text{ dla }k,i = 1, \ldots, n 
\end{equation}
Wtedy: 
\begin{equation}\label{eq:std_inner}
    \inner{\textbf{x}}{\textbf{y}} = \sum_{i,k=1}^{n}{h_{ik} x_{i} \bar{y}_{k}}
\end{equation}
W szczególności:

\begin{equation}
    \inner{\textbf{x}}{\textbf{x}} = \sum_{i,k=1}^{n}{h_{ik} x_{i} \bar{x}_{k}}
\end{equation}
gdzie z \eqref{eq:def_h} oraz własności (1), wynika, że $h_{ik} = \bar{h}_{ki}$.

% TODO: Potrzebujemy to dodać formę hermitowską?

W ten sposób zdefiniowaliśmy iloczyn skalarny dla bazy $\textbf{e}_1, \ldots, \textbf{e}_n$. Ponadto jeżeli baza $\textbf{e}_1, \ldots, \textbf{e}_n$ jest ortonormalna, to znaczy taka, że każdy wektor bazy jest ortogonalny do pozostałych oraz ma normę jeden:
\begin{equation}\label{eq:def_ortonormal_base}
    \inner{\textbf{e}_i}{\textbf{e}_k} = \delta_{ik} = \begin{cases}
    0, \text{ dla } i \neq k \\
    1, \text{ dla } i = k
\end{cases}
\end{equation}
Wtedy dla bazy ortonormalnej iloczyn skalarny przyjmuje postać, którą będziemy nazywać \textbf{standardową}:
\begin{equation}\label{eq:std_normal_inner}
    \inner{\textbf{x}}{\textbf{y}} = \sum_{i=1}^{n}{x_{i} \bar{y}_{i}}
\end{equation}
\begin{equation}
    \inner{\textbf{x}}{\textbf{x}} = \sum_{i=1}^{n}{x_{i} \bar{x}_{i}} = \sum_{i=1}^{n}{|x_{i}|^2}
\end{equation}
Dla wektorów \textbf{x}, \textbf{y} ze współrzędnymi $x_i$ i $y_i$ dla $i = 1, \ldots, n$ w bazie $\textbf{e}_1, \ldots, \textbf{e}_n$
przestrzeni euklidesowej iloczyn skalarny redukuje się do:
\begin{equation}\label{eq:std_euclid_inner}
    \inner{\textbf{x}}{\textbf{y}} = \sum_{k,i=1}^{n}{h_{ik} x_{i} {y}_{k}}
\end{equation}
Dla bazy ortonormalnej:
\begin{equation}\label{eq:std_euclid_normal_inner}
    \inner{\textbf{x}}{\textbf{y}} = \sum_{i=1}^{n}{x_{i} {y}_{i}}
\end{equation}
\begin{equation}
    \inner{\textbf{x}}{\textbf{x}} = \sum_{i=1}^{n}{x_{i}^2}
\end{equation}

\

\begin{definition}
Odwzorowanie $A \colon X \to Y $, gdzie $X$ i $Y$ są przestrzeniami liniowymi nad ciałem liczb $\K$, nazywa się \textbf{operatorem liniowym} jeżeli spełnia następujące warunki dla dowolnych dwóch wektorów $x, y \in X$ oraz dowolnego skalara $a \in \K$:
\begin{enumerate}[leftmargin=1cm, itemsep=2pt, label=(\arabic*)]
    \item $A(x + y) = A(x) + A(y)$ (zachowuje dodawanie wektorów)
    \item $A(ax) = aA(x)$ (zachowuje mnożenie przez skalar)
\end{enumerate}
\end{definition}

Jeśli przestrzenie liniowe $X$ i $Y$ mają skończony wymiar $\dim X = n$, $\dim Y = m$, oraz ustalimy w $X$ bazę $\mathbf{e}_1, \ldots, \mathbf{e}_n$, a w $Y$ bazę $\mathbf{g}_1, \ldots, \mathbf{g}_m$, to odwzorowaniu liniowemu $A$ odpowiada macierz $m \times n$ w następującym sensie: jeżeli wektorowi $\mathbf{x} = \sum_{i=1}^{n} x_i \mathbf{e}_i \in X$ operator $A$ przyporządkowuje wektor $\mathbf{y} = \sum_{i=1}^{m} y_i \mathbf{g}_i \in Y$, to współrzędne tych wektorów są związane wzorem $\mathbf{y} = A\mathbf{x}.$

Uwaga: W dalszej części pracy, ustaliwszy bazę w przestrzeni $X$ i $Y$, będziemy utożsamiać operator liniowy $A \colon X \to Y$ z jego macierzą $A$ w tej bazie, pisząc po prostu $A$ zarówno dla operatora, jak i dla odpowiadającej mu macierzy.

\section{Liniowe operatory w przestrzeni euklidesowej}

Rozpatrzmy $n$-wymiarową przestrzeń euklidesową $V$. Niech będzie dany dowolny operator liniowy $A$ w tej przestrzeni.

\begin{definition}
    Operator liniowy $A^T$ nazywa się \textbf{operatorem transponowanym do operatora} $A$, jeśli dla dowolnych dwóch wektorów $x, y \in V$:
    \begin{equation}\label{eq:def_transposed_operator}
        \inner{Ax}{y} = \inner{x}{A^Ty}
    \end{equation}
\end{definition}
\begin{lemma}\label{lemma:existing_of_ortonormal_base}
W każdej skończonie wymiarowej przestrzeni unitarnej $V$ można wybrać bazę ortonormalną.
\end{lemma}   
\begin{proof}
    Dowolną bazę $V$ można zortogonalizować metodą Grama–Schmidta \cite[s. 140--142]{mostowski1977}, a następnie znormalizować otrzymane wektory, co daje bazę ortonormalną.
\end{proof}

\begin{theorem}
Dla każdego liniowego operatora $A$ w przestrzeni euklidesowej $V$ zawsze istnieje operator transponowany $A^T$ i to tylko jeden.
\end{theorem}
\begin{proof}
\textbf{Istnienie.} Niech $\mathbf{e}_1, \ldots, \mathbf{e}_n$ jest bazą ortonormalną przestrzeni $V$. Dowolny wektor $y \in V$ w bazie ortonormalnej da się zapisać w postaci:
$$ y = \sum_{k=1}^{n}{\inner{y}{\mathbf{e}_k}\mathbf{e}_k} $$
    Wtedy dla wektora $y$ i szukanego operatora $A^T$:
\begin{equation}
        A^Ty = \sum_{k=1}^{n}{\inner{A^Ty}{\mathbf{e}_k}\mathbf{e}_k}
\end{equation}
    Na mocy Lematu \ref{lemma:existing_of_ortonormal_base}:
\begin{equation}
        A^Ty = \sum_{k=1}^{n}{\inner{y}{A\mathbf{e}_k}\mathbf{e}_k}
\end{equation}

Operator ten jest liniowy, ponieważ iloczyn skalarny jest liniowy.
Pokażemy, że tak zdefiniowany operator spełnia warunek \eqref{eq:def_transposed_operator}. Niech $x, y \in V$ będą dowolne. Zapiszmy $x$ w bazie ortonormalnej: $x = \sum_{k=1}^n \inner{x}{\mathbf{e}_k}\mathbf{e}_k$. Wówczas
\begin{equation}
    \inner{x}{A^Ty} = \Big\langle \sum_{k=1}^n \inner{x}{\mathbf{e}_k}\mathbf{e}_k,\ \sum_{i=1}^n \inner{y}{A\mathbf{e}_i}\mathbf{e}_i \Big\rangle = \sum_{k=1}^n \inner{x}{\mathbf{e}_k}\inner{y}{A\mathbf{e}_k},
\end{equation}
gdzie skorzystaliśmy z ortonormalności $\base{e}$. 
Z drugiej strony
\begin{equation}
    \inner{Ax}{y} = \inner{
        A\Big(\sum_{k=1}^n \inner{x}{\mathbf{e}_k}\mathbf{e}_k\Big)
    }{y} = \sum_{k=1}^n \inner{x}{\mathbf{e}_k}\inner{A\mathbf{e}_k}{y} = \sum_{k=1}^n \inner{x}{\mathbf{e}_k}\inner{y}{A\mathbf{e}_k},
\end{equation}
gdzie ostatnia równość wynika z przemienności iloczynu skalarnego w przestrzeni euklidesowej. 
Prawe strony obu równości są identyczne, więc $\inner{Ax}{y} = \inner{x}{A^Ty}$ dla wszystkich $x,y \in V$, co dowodzi istnienia $A^T$.

\textbf{Jednoznaczność.} Załóżmy, że operatory $B_1, B_2$ oba spełniają warunek \eqref{eq:def_transposed_operator}, tzn. $\inner{Ax}{y} = \inner{x}{B_1y} = \inner{x}{B_2y}$ dla wszystkich $x,y \in V$. Wtedy dla dowolnych $x, y \in V$:
\begin{equation}
    \inner{x}{B_1y - B_2y} = 0.
\end{equation}
Ponieważ $x \in V$ było dowolne, a iloczyn skalarny w $V$ jest niezdegenerowany (tzn. $\inner{x}{z}=0$ dla wszystkich $x$ różnych od zera implikuje $z=0$), otrzymujemy $B_1y - B_2y = 0$ dla każdego $y \in V$, czyli $B_1 = B_2$.
\end{proof}

\begin{theorem}
    W przestrzeni euklidesowej $V$ z bazą ortonormalną, transponowanemu operatorowi $A^T$ odpowiada macierz transponowana.
\end{theorem}
\begin{proof}
    Niech $A$ będzie operatorem liniowym w przestrzeni euklidesowej $V$ z  ortonormalną bazą $\base{e}$ (Lemat \ref{lemma:existing_of_ortonormal_base}) któremu odpowiada macierz $(a_{ik})$, gdzie $i, k = 1, \ldots, n$.

    Zastosujmy do dowolnego wektora $\mathbf{e}_k$ z bazy $\base{e}$ operator liniowy $A$:
    \begin{equation}\label{eq:local_Ae_def}
        A\mathbf{e}_k = \sum_{i=1}^{n}{a_{ik}\mathbf{e}_i} = a_{1k}\mathbf{e}_1 + \ldots + a_{ik}\mathbf{e}_i + \ldots + a_{nk}\mathbf{e}_n
    \end{equation}
    Pomnożymy skalarnie obie strony \eqref{eq:local_Ae_def} przez dowolny wektor bazowy $\mathbf{e}_i$, gdzie $i \in \{1, \ldots, n\}$:
    $$\inner{A\mathbf{e}_k}{\mathbf{e}_i} = \inner{
        a_{1k}\mathbf{e}_1 + \ldots + a_{ik}\mathbf{e}_i + \ldots + a_{nk}\mathbf{e}_n
    }{\mathbf{e}_i}$$
    $$\inner{A\mathbf{e}_k}{\mathbf{e}_i} = a_{1k}\inner{\mathbf{e}_1}{\mathbf{e}_i} + \ldots + a_{ik}\inner{\mathbf{e}_i}{\mathbf{e}_i} + \ldots + a_{nk}\inner{\mathbf{e}_n}{\mathbf{e}_i}$$
    Z definicji bazy ortonormalnej \eqref{eq:def_ortonormal_base}:
    $$\inner{A\mathbf{e}_k}{\mathbf{e}_i} = a_{1k} \times 0 + \ldots + a_{ik}\times 1 + \ldots + a_{nk} \times 0$$
    Więc:
    \begin{equation}\label{eq:def_ortonormal_coeffs}
        a_{ik} = \inner{A\mathbf{e}_k}{\mathbf{e}_i}
    \end{equation}
    Wzór \eqref{eq:def_ortonormal_coeffs} jest wzorem na współczynniki macierzy operatora liniowego w bazie ortonormalnej.
    %TODO: We need to split that formula and show this before theorem about existing transposed op%
    Niech w tej samej bazie operatorowi $A^T$ odpowiada macierz $(a^T_{ik})$, gdzie $i, k = 1, \ldots, n$.
    Wtedy ze wzoru \eqref{eq:def_ortonormal_coeffs}:
    \begin{equation}\label{eq:local_def_a_t_el}
        a^T_{ik} = \inner{A^T\mathbf{e}_k}{\mathbf{e}_i}
    \end{equation}
    Z definicji operatora transponowanego \eqref{eq:def_transposed_operator}:
    $$\inner{A^T\mathbf{e}_k}{\mathbf{e}_i} = \inner{\mathbf{e}_k}{A\mathbf{e}_i}$$
    \begin{equation}\label{eq:local_inner_t_eq_to_inner}
        \inner{A^T\mathbf{e}_k}{\mathbf{e}_i} = \inner{A\mathbf{e}_i}{\mathbf{e}_k}
    \end{equation}
    Z \eqref{eq:def_ortonormal_coeffs}, \eqref{eq:local_def_a_t_el} oraz \eqref{eq:local_inner_t_eq_to_inner} wynika, że:
    $$a^T_{ik} = a_{ki}, \quad (i, k = 1, \ldots, n)$$
    A zatem operatorowi transponowanemu $A^T$ odpowiada macierz transponowana.
\end{proof}

\begin{theorem}\label{theorem:properties_of_transposed_operator}
    Z definicji operatora transponowanego mamy następujące własności:
    \begin{enumerate}[leftmargin=1cm, itemsep=2pt, label=(\arabic*)]
        \item $(A^T)^T = A$
        \item $(A + B)^T = A^T + B^T$
        \item $(\lambda A)^T = \lambda(A^T)$, gdzie $\lambda \in \R$
        \item $(AB)^T = B^TA^T$
    \end{enumerate}
\end{theorem}
\begin{proof}
Dla (1): $(A^T)^T = (a^T_{ik})^T = (a_{ki})^T = (a_{ik}) = A$, gdzie $i, k \in \N$

Dla (2):
$(A + B)^T = (a_{ik} + b_{ik})^T = a_{ki} + b_{ki} = a^T_{ik} + b^T_{ik} = A^T + B^T$, gdzie $i, k \in \N$

Dla (3):
$(\lambda A)^T = (\lambda a^T_{ik}) = \lambda(a^T_{ik}) = \lambda(A^T)$, gdzie $\lambda \in \R, i, k \in \N$

Dla (4):
Zakładamy, że możemy wymnożyć macierze $A$ i $B$, a więc liczba kolumn macierzy $A_{(p \times n)}$ zgadza się z liczbą wierszy macierzy $B_{(n \times m)}$. Z tego wynika, że liczba wierszy macierzy $A^T_{(n \times p)}$, zgadza się z liczbą kolumn macierzy $B^T_{(m \times n)}$, gdzie $n, p, m \in \N$.
Element na pozycji $(i, k)$ macierzy $(AB)$ z defnicji mnożenia dany jest wzorem: 

$$(AB)_{(i, k)} = \sum_{r=1}^{n}{a_{ir}b_{rk}}$$
Element na pozycji $(i, k)$ macierzy $(AB)^T$ dany jest wzorem:
$$(AB)^T_{(i, k)} = (AB)_{(k, i)} = \sum_{r=1}^{n}{a_{kr}b_{ri}}$$
Dla macierzy $(B^TA^T)$ element $(i, k)$ z definicji mnożenia macierzy jest dany wzorem:
$$(B^TA^T)_{(i, k)} = \sum_{r=1}^{n}{b^T_{ir}a^T_{rk}} = \sum_{r=1}^{n}{b_{ri}a_{kr}} = \sum_{r=1}^{n}{a_{kr}b_{ri}}$$
Z tego wynika, że $(B^TA^T)_{(i, k)} = (AB)^T_{(i, k)}$, więc $(AB)^T = B^TA^T$.

\end{proof}

\begin{definition}\label{def:normal_operator}
Operator liniowy $A$ nazywa się \textit{normalnym}, jeśli:
$$AA^T = A^TA$$
\end{definition}
\begin{theorem}
    Operator normalny $A$ w przestrzeni euklidesowej $V$ z bazą ortonormalną ma macierz normalną.
\end{theorem}
\begin{proof}

\end{proof}

\begin{definition}\label{def:symmetric_operator}
Operator liniowy $S$ nazywa się \textit{symetrycznym}, jeśli:
$$S = S^T$$
\end{definition}
\begin{theorem}\label{theorem:sym_op_has_same_matrix}
    Operator symetryczny $S$ w przestrzeni euklidesowej $V$ z bazą ortonormalną ma macierz symetryczną.
\end{theorem}
\begin{proof}
    Z definicji operatora symetrycznego dla dowolnych dwóch wektorów $\textbf{e}_i$, $\textbf{e}_k$ z bazy ortonormalnej $\base{e}$, gdzie $i, k \in \N$:
    $$\inner{S\textbf{e}_i}{\textbf{e}_k} = \inner{\textbf{e}_i}{S\textbf{e}_k}$$
    Korzystając ze wzoru \eqref{eq:def_ortonormal_coeffs} oraz przemienności iloczynu skalarnego:
    $$ a_{ik} = \inner{S\textbf{e}_i}{\textbf{e}_k} = \inner{\textbf{e}_i}{S\textbf{e}_k} = \inner{S\textbf{e}_k}{\textbf{e}_i} = a_{ki}, $$ gdzie $i, k \in \{1, \dots, n\}$. Z tego wynika, że operatorowi symetrycznemu odpowiada macierz symetryczna.
\end{proof}

\begin{definition}\label{def:skew_symmetric_operator}
Operator liniowy $K$ nazywa się \textit{antysymetrycznym}, jeśli:
$$K = -K^T$$
\end{definition}

\begin{theorem}
    Operator antysymetryczny $K$ w przestrzeni euklidesowej $V$ z bazą ortonormalną ma macierz antysymetryczną.
\end{theorem}
\begin{proof}
    Dowód jest podobny do dowodu Twierdzenia \ref{theorem:sym_op_has_same_matrix}.

    Z definicji operatora antysymetrycznego dla dowolnych dwóch wektorów $\textbf{e}_i$, $\textbf{e}_k$ z bazy ortonormalnej $\base{e}$, gdzie $i, k \in \N$:
    $$\inner{K\textbf{e}_i}{\textbf{e}_k} = \inner{\textbf{e}_i}{-K\textbf{e}_k}$$
    Korzystając ze wzoru \eqref{eq:def_ortonormal_coeffs} oraz z jednorodności i przemienności iloczynu skalarnego:
    $$ a_{ik} = \inner{K\textbf{e}_i}{\textbf{e}_k} = \inner{\textbf{e}_i}{-K\textbf{e}_k} = \inner{-K\textbf{e}_k}{\textbf{e}_i} = -\inner{K\textbf{e}_k}{\textbf{e}_i} = -a_{ki}, $$ gdzie $i, k \in \{1, \dots, n\}$. Z tego wynika, że operatorowi antysymetrycznemu odpowiada macierz antysymetryczna.
\end{proof}

\begin{definition}\label{def:ortogonal_operator}
    Operator liniowy $O$ nazywa się \textit{ortogonalnym} lub inaczej \textit{izometrią liniową}, jeśli zachowuje metrykę przestrzeni, tj. jeśli dla dowolnych dwóch wektorów $x, y \in V$:
    \begin{equation}\label{eq:def_ortogonal_operator}
        \inner{Ox}{Oy} = \inner{x}{y}
    \end{equation}
\end{definition}
Równość \eqref{eq:def_ortogonal_operator} możemy przepisać w następujący sposób na mocy \eqref{eq:def_transposed_operator}:
$$\inner{Oa}{b} = \inner{a}{O^Tb}$$
Podstawiając pod $a = x$, a pod $b = Oy$ otrzymujemy:
$$\inner{Ox}{Oy} = \inner{x}{O^TOy} = \inner{x}{y}$$
Z tego wynika:
\begin{equation}\label{eq:def_ortogonal_operator_2}
    O^TO = I,
\end{equation} gdzie $I$ macierz jednostkowa. 

\begin{theorem}
    Operator ortogonalny $O$ w przestrzeni euklidesowej $V$ z bazą ortonormalną ma macierz ortogonalną.
\end{theorem}
\begin{proof}
    Załóżmy, że $O$ jest operatorem ortogonalnym w przestrzeni euklidesowej z bazą ortonormalną $\base{e}$, i odpowiednią mu macierzą, oznaczymy ją jako $A$.
    Ponieważ $\base{e}$ jest bazą ortonormalną, to $O(\mathbf{e}_1), \ldots, O(\mathbf{e}_n)$ też jest bazą ortonormalną, dlatego że $O$ zachowuje iloczyn skalarny tych wektorów. Te wektory traktowane jako wektory kolumnowe tworzą macierz ortogonalną, tworząc jednocześnie macierz $A$.

    Na odwrót, załóżmy, że macierz $A$ jest ortogonalna. Niech $x, y \in V$. Wówczas, korzystając z tw. \ref{theorem:properties_of_transposed_operator} (4 własność):
    $$\inner{O(x)}{O(y)} = \inner{Ax}{Ay} = (Ax)^TAy = x^TA^TAy = x^Ty = \inner{x}{y},$$
    a to oznacza, że $O$ jest izometrią liniową.
\end{proof}

Operatory symetryczne, antysymetryczne oraz ortogonalne są szczególnymi przypadkami operatora normalnego.
Jak widzimy z powyższego, utożsamiamy macierze symetryczne, antysymetryczne, ortogonalne z operatorami liniowymi lub inaczej, w przestrzeni euklidesowej każda z tych macierzy reprezentuje przekształcenie liniowe.

\section{Twierdzenie spektralne}

\section{Postać normalna macierzy rzeczywistych symetrycznych, antysymetrycznych oraz ortogonalnych}

W tym podrozdziale omówimy postaci normalne dla macierzy rzeczywistych symetrycznych, antysymetrycznych oraz ortogonalnych i pokażemy kilka przykładów. Także omówimy dlaczego przypadek zespolony jest bardziej skomplikowany.




%  Wzór na współczynniki macierzy w bazie ortonormalnej.
% TODO:
% 4. Ostatni podrozdział: forme normalne i przeprowadzić dowód (ale bez tw spektralnego) 
% Skorzystać z tego:
% W dowodach poniższych twierdzeń nie będziemy korzystać z ogólnego twierdzenia spektralnego dla operatorów normalnych, lecz przeprowadzimy dowody bezpośrednio, wykorzystując jedynie własności iloczynu skalarnego oraz definicje odpowiednich klas operatorów.
% 5. Opisać dla czego przypadek zespolony jest cięższy i wprowdzić rozdział 2 (Dzielniki elementarne oraz postać Jordana)

% Udowowdnić Tw Spektralne dla operatora normalnego, lepiej pokazać że operator normalny ma baze ortonormalną
% Pokazać że macierz w R ma wartości własne sprzężone(parami)
% Pokazać że operator liniowy w R 
